\chapter{Split-Signature Clifford Geometry}
\label{ch:ch01_clifford}

This chapter fixes the algebra and makes no physical reading, with two forward references: \cref{r3:prop:clock} uses the Fisher--Rao clock
of \cref{ch:ch02_lorentz}, and \eqref{r3:eq:Xfock} writes a vector in the electric and magnetic components of \cref{ch:ch03_readings}. We work
with the real Clifford algebra $\Cl(3,3)\cong M_8(\R)$ of $\R^{3,3}$ in an explicit realisation, with the blade rules and the signature rule,
which separates compact bivectors from boosts. The trace form on an orbit splits by the Cartan involution (\cref{thm:ghost}); for actions built
on this trace metric, this decides the sign of the kinetic terms later on. The
spinor module is then described as the forms on a null three-plane, with its octants and Fock grading, and the tower $\Cl(3n,3n)$ is
defined. Every statement is algebra, proved in the text and checked exactly (\cref{app:manifest}). \Cref{ch:ch02_lorentz} introduces
the Lorentz structures and the clock; \cref{ch:ch03_readings} the two readings of a vector of the arena.


% --- source: tex/s02_arena#0 ((head))


% --- source: tex/s02_arena#1 (The split space and its Clifford algebra)
\section{The split space and its Clifford algebra}
\label{sec:arena}

% --- source: tex/s02_arena#2 (A fixed realisation)
\subsection{A fixed realisation}

Let $V=\R^{3,3}$ with orthonormal basis $e_1,\dots,e_6$, $e_a^2=\eta_{aa}$ and $\eta=\operatorname{diag}(1,1,1,-1,-1,-1)$.
We call $V^+=\langle e_1,e_2,e_3\rangle$ the temporal triple and $V^-=\langle e_4,e_5,e_6\rangle$ the spatial triple. The
Clifford algebra $\Cl(3,3)$ is generated by the $e_a$ subject to $e_ae_b+e_be_a=2\eta_{ab}\id$; it is isomorphic to
$M_8(\R)$, and we fix once and for all its realisation on $S=\R^8=(\R^2)^{\otimes3}$,
\begin{equation}
\label{eq:realisation}
\begin{aligned}
e_1&=\sigma_x\otimes\id\otimes\id, & e_2&=\sigma_z\otimes\sigma_x\otimes\id, & e_3&=\sigma_z\otimes\sigma_z\otimes\sigma_x,\\
e_4&=\varepsilon\otimes\id\otimes\id, & e_5&=\sigma_z\otimes\varepsilon\otimes\id, & e_6&=\sigma_z\otimes\sigma_z\otimes\varepsilon,
\end{aligned}
\end{equation}
with $\sigma_x=\left(\begin{smallmatrix}0&1\\1&0\end{smallmatrix}\right)$,
$\sigma_z=\left(\begin{smallmatrix}1&0\\0&-1\end{smallmatrix}\right)$ and
$\varepsilon=\left(\begin{smallmatrix}0&-1\\1&0\end{smallmatrix}\right)$. All six matrices are signed permutation
matrices. For a subset $I=\{a_1<\dots<a_r\}$ we write $e_I=e_{a_1}\cdots e_{a_r}$ for the corresponding coordinate blade of grade $r$;
the $64$ blades form a basis of $M_8(\R)$.

The tower is built by graded tensor products. At $n$ constituents the spinor module is $S_n=(\R^8)^{\otimes n}$ and the
generators of the $k$-th constituent are
\begin{equation}
\label{eq:tower}
e^{(k)}_a=\underbrace{\omega\otimes\cdots\otimes\omega}_{k-1}\otimes e_a\otimes\id\otimes\cdots\otimes\id,
\qquad \omega:=e_1e_2\cdots e_6 .
\end{equation}
Since $\omega^2=\id$ and $\omega$ anticommutes with every $e_a$, generators of different constituents anticommute, and the
$6n$ matrices $e^{(k)}_a$ generate $\Cl(3n,3n)\cong M_{8^n}(\R)$. When we need a single index we write
$E_{6(k-1)+a}=e^{(k)}_a$, so that at three constituents the eighteen generators are $E_1,\dots,E_{18}$ and the temporal
ones are $E_A$ with $A\in\{1,2,3,7,8,9,13,14,15\}$.

Two elementary rules carry most of the finite-dimensional arguments of this paper.

\begin{lemma}[Blade rules]
\label{lem:blade}
For coordinate blades $e_I,e_J$ of grades $|I|,|J|$,
\begin{equation}
\label{eq:bladecomm}
e_Ie_J=(-1)^{|I||J|-|I\cap J|}\,e_Je_I ,
\end{equation}
and for a coordinate blade $B$ of grade $r$
\begin{equation}
\label{eq:transpose}
B^{\!\top}=(-1)^{r(r-1)/2}\Bigl(\prod_{a\in I}\eta_{aa}\Bigr)B=B^2B .
\end{equation}
In particular a coordinate blade is antisymmetric if and only if it squares to $-\id$, and symmetric if and only if it squares to $+\id$.
\end{lemma}

\begin{proof}
For \eqref{eq:bladecomm}, move each generator of $e_J$ past the generators of $e_I$: it anticommutes with the $|I|$ of them
it differs from and commutes with itself when it occurs in $I$, which gives the exponent. For \eqref{eq:transpose}, the
matrices $\sigma_x,\sigma_z$ are symmetric and $\varepsilon$ is antisymmetric, so each $e_a$ satisfies $e_a^{\!\top}=\eta_{aa}e_a$;
transposing a product reverses it, and restoring the order costs $(-1)^{r(r-1)/2}$. The same two factors give
$B^2=(-1)^{r(r-1)/2}\prod_{a\in I}\eta_{aa}\,\id$, so $B^{\!\top}=B^2B$.
\end{proof}

\begin{example}
\label{ex:blades}
The bivector $e_2e_3$ has $(e_2e_3)^2=-e_2^2e_3^2=-\id$, so it is antisymmetric and $\exp(\theta e_2e_3)=\cos\theta\,\id+
\sin\theta\,e_2e_3$ is a rotation of period $2\pi$. The bivector $e_1e_4$ has $(e_1e_4)^2=-e_1^2e_4^2=+\id$, so it is
symmetric and $\exp(\theta e_1e_4)=\cosh\theta\,\id+\sinh\theta\,e_1e_4$ is a boost, unbounded in $\theta$. These two
bivectors share no index, $|I||J|-|I\cap J|=4-0=4$ is even, and they commute; so do $e_2e_3$ and $e_5e_6$. The boost $e_1e_4$ and the
rotation $e_1e_2$ share the index $1$, $|I||J|-|I\cap J|=4-1=3$ is odd, and they anticommute.
\end{example}

% --- source: tex/s02_arena#3 (The signature rule)
\subsection{The signature rule}
\label{sec:sigrule}

\begin{theorem}[The signature rule for coordinate bivectors]
\label{thm:sig}
For $a\neq b$, $(e_ae_b)^2=-e_a^2e_b^2$. A bivector therefore squares to $-\id$ when its two directions have the same
sign and to $+\id$ when they have opposite signs, and by \cref{lem:blade} the first kind are exactly the antisymmetric
bivectors. The three time--time and three space--space bivectors span the maximal compact subalgebra
$\so3\oplus\so3\subset\so{3,3}$, and the nine mixed bivectors are boosts.
\end{theorem}

\begin{corollary}[A mixed bivector generates no compact $U(1)$]
\label{cor:noboost}
For a mixed bivector $B$, $\exp(\theta B)=\cosh\theta\,\id+\sinh\theta\,B$ generates $\R$ and not $U(1)$: it carries no
phase and defines no compact charge.
\end{corollary}

The obstruction to a boost as a gauge generator lies in its action on the module, not in the kinetic metric of the gauge field. For an
abelian algebra that metric decides nothing: on $\R B$ every scalar product is invariant, and $A_\mu=a_\mu B$ has the ordinary, healthy
abelian kinetic term. What fails is unitarity. No positive-definite $H$ on $S$ satisfies $B^{\!\top}H+HB=0$, since an eigenvector $Bv=\pm v$
would give $v^{\!\top}Hv=0$. A boost is therefore not unitary for any fixed positive form on the module, and the gauge form $-\Tr$ has the
wrong sign on it. Relative to the chosen positive matter metric $H$, we call an internal generator $X$ \emph{gauge-admissible} when
$X^{\!\top}H+HX=0$; in the orthonormal realisation used here, $H=\id$, this is $X^{\!\top}=-X$, that is, $X$ is compact. This is the sense in which the corollary is used later (\cref{e1:thm:Nrule}, \cref{ch:ch09_bfield,ch:ch13_compatibility}),
and the conclusions drawn there hold under it. Allowing an indefinite invariant matter form removes this positive-metric obstruction.
Whether physical ghosts then occur depends on the chosen action and its constraint reduction (an invariant constraint can remove the
negative direction); an unconstrained quadratic kinetic term with propagating negative-sign directions is unacceptable. The present model
therefore imposes a fixed positive matter metric and restricts its internal gauge generators accordingly.

The criterion that matters for a gauge generator is antisymmetry with respect to the positive form on the module that a kinetic term
supplies, equivalently membership of the $+1$ eigenspace of the Cartan involution $\theta(X)=-X^{\!\top}$
\cite{Knapp,Helgason}. For a single blade this coincides with $X^2=-\id$ by \eqref{eq:transpose}, which is why the rule is
decisive for bivectors. For general elements it is not: $\operatorname{diag}(J_2,2J_2)$, with $J_2$ the standard $2\times2$
rotation generator, generates a group of period $2\pi$ although its square is $\operatorname{diag}(-\id,-4\id)$.

\begin{figure}[t]
\centering
\begin{tikzpicture}[scale=0.82]
\foreach \i in {1,...,6}{
  \node at (\i-0.5,0.35) {\small $e_{\i}$};
  \node at (-0.35,-\i+0.5) {\small $e_{\i}$};
}
\foreach \i in {1,...,6}{
  \foreach \j in {1,...,6}{
    \pgfmathtruncatemacro{\ti}{\i<4 ? 1 : 0}
    \pgfmathtruncatemacro{\tj}{\j<4 ? 1 : 0}
    \ifnum\j>\i
      \ifnum\ti=\tj
        \fill[tcol!35] (\j-1,-\i+1) rectangle (\j,-\i);
      \else
        \fill[scol!25] (\j-1,-\i+1) rectangle (\j,-\i);
      \fi
    \fi
    \ifnum\j=\i \fill[lgray] (\j-1,-\i+1) rectangle (\j,-\i); \fi
  }
}
\draw[step=1,mgray,thin] (0,-6) grid (6,0);
\draw[very thick] (0,0) rectangle (6,-6);
\draw[thick] (3,0) -- (3,-6); \draw[thick] (0,-3) -- (6,-3);
% pole operators
\foreach \i in {1,2,3}{\pgfmathtruncatemacro{\jj}{\i+3}\node at (\jj-0.5,-\i+0.5) {\footnotesize $\omega_{\i}$};}
\node at (2.5,-1.5) {\footnotesize $J$};
\fill[tcol!35] (6.6,-1.05) rectangle ++(0.32,0.32); \draw[mgray] (6.6,-1.05) rectangle ++(0.32,0.32);
\node[align=left,anchor=west] at (7.05,-0.9) {\small same sign: $B^2=-\id$, $B^{\!\top}=-B$,};
\node[align=left,anchor=west] at (7.05,-1.5) {\small compact; $3+3$ generators of $\so3\oplus\so3$};
\fill[scol!25] (6.6,-2.65) rectangle ++(0.32,0.32); \draw[mgray] (6.6,-2.65) rectangle ++(0.32,0.32);
\node[align=left,anchor=west] at (7.05,-2.5) {\small mixed sign: $B^2=+\id$, $B^{\!\top}=B$,};
\node[align=left,anchor=west] at (7.05,-3.1) {\small boosts; $9$ generators};
\node[align=left,anchor=west] at (6.6,-4.1) {\small $\omega_i=e_ie_{i+3}$: the pole operators,};
\node[align=left,anchor=west] at (6.6,-4.7) {\small \phantom{$\omega_i$} commuting involutions};
\node[align=left,anchor=west] at (6.6,-5.5) {\small $J=e_2e_3$: the interior charge};
\end{tikzpicture}
\caption{The fifteen bivectors $e_ae_b$, $a<b$, of $\Cl(3,3)$. Blue cells pair two directions of the same sign and generate
compact rotations; red cells pair a temporal with a spatial direction and generate boosts (\cref{thm:sig}). The diagonal of
the red block holds the three pole operators, which grade the algebra (\cref{prop:grading}); the cell $(2,3)$ holds the
bivector $J$, the compact rotation that commutes with the interior Lorentz algebra (\cref{prop:6plus8plus1}) and is used as the imaginary unit
in \cref{ch:ch02_lorentz}. It is not identified with electromagnetism: the electromagnetic generator is constructed in \cref{ch:ch08_photon}.}
\label{fig:bivectors}
\end{figure}

% --- source: tex/s02_arena#4 (The pole operators and a multigrading)
\subsection{The pole operators and a multigrading}
\label{sec:grading}

The mixed bivectors are not idle. The three \emph{pole operators}
\[
\omega_i:=e_ie_{i+3}\qquad(i=1,2,3)
\]
are boosts, so $\omega_i^2=+\id$, and they commute pairwise by \eqref{eq:bladecomm}, since any two share no index. As operators on $S$ the
two kinds of bivector differ: one of square $-\id$ is a complex structure, while each $\omega_i$ has real $\pm1$ eigenspaces, which grade the
module. By conjugation on the algebra either kind induces an involution: conjugation by $J=e_2e_3$, of square $-\id$, has $\pm1$ eigenspaces of
dimension $32$ each.

\begin{proposition}[A $\Z_2^3$ grading of $\Cl(3,3)$]
\label{prop:grading}
For a blade $e_I$ put $\chi_i(e_I)=(-1)^{|I\cap\{i,i+3\}|}$. Then $\omega_ie_I\omega_i^{-1}=\chi_i(e_I)\,e_I$, and
$\Cl(3,3)=\bigoplus_{\sigma\in\{\pm\}^3}\Cl_\sigma$ is a multiplicative grading with eight sectors of dimension eight.
\begin{enumerate}[label=(\roman*),leftmargin=2.2em]
\item The joint spectrum of $(\omega_1,\omega_2,\omega_3)$ on $S$ consists of all eight sign triples, each on a line, so the
pole operators cut $\R^8$ into eight one-dimensional eigenlines.
\item The vectors occupy the three weights with exactly one minus sign, two each ($e_i$ and $e_{i+3}$); the bivectors occupy
the four even weights with multiplicities $3,4,4,4$; the all-minus weight is spanned by eight trivectors, namely $e_1e_2e_3$,
$e_4e_5e_6$ and six mixed ones.
\item Correspondingly $\so{3,3}=\mathfrak g_0\oplus\mathfrak g_1\oplus\mathfrak g_2\oplus\mathfrak g_3$ is graded by the Klein
four-group, of dimensions $3+4+4+4$, with $\mathfrak g_0=\operatorname{span}\{\omega_1,\omega_2,\omega_3\}$ abelian,
self-normalising and consisting of semisimple elements with real eigenvalues: a maximally split Cartan subalgebra of
$\so{3,3}\cong\slr{4,\R}$.
\end{enumerate}
\end{proposition}

\begin{proof}
The conjugation sign is \eqref{eq:bladecomm} with $|J|=2$: $(-1)^{2|I|-|I\cap\{i,i+3\}|}$. The weight is multiplicative because
$|I\triangle I'|$ and $|I|+|I'|$ have the same parity in each pair. The census in (ii) is a count of how many generators of
each pair a blade contains. For (iii), a bivector inside one pair is a pole operator, of weight $(+,+,+)$; a bivector with one
generator in each of two pairs has two minus signs, and each of the three such weights holds $2\times2=4$ bivectors. The
pole operators commute, are symmetric with eigenvalues $\pm1$, and the centraliser of $\mathfrak g_0$ in $\so{3,3}$ is
$\mathfrak g_0$ itself, because every other bivector has a nontrivial weight; its dimension equals the real rank three of
$\so{3,3}$.
\end{proof}

The three characters are grading weights, not metric signatures: $e_1$ and $e_4$ carry the same weight and opposite squares,
so nothing in the grading tells a temporal from a spatial direction. The decomposition in (iii) is the root decomposition of
$\slr{4,\R}$ with respect to a split Cartan subalgebra, grouped by the coordinate parity of the roots, $Q(D_3)/2\Z^3$, a group of order four
(not the root lattice modulo twice itself, $Q(D_3)/2Q(D_3)$, which has order eight); gradings of Clifford and
classical Lie algebras by $\Z_2^n$ are well known \cite{MosnaMirallesVaz}, and we use this one because the pole operators will
reappear in the field-strength reading as the three paired planes $(e_i,e_{i+3})$, $i=1,2,3$, that carry one complex component each.

\begin{example}[The paired-mode projector]
\label{ex:pairedmode}
Let $\varpi=\omega_1+\omega_2+\omega_3$ be the sum of the pole operators and $\omega=e_1e_2\cdots e_6$ the volume element of
\eqref{eq:tower}. On the joint eigenline of signs $(\varepsilon_1,\varepsilon_2,\varepsilon_3)$, $\varpi=\varepsilon_1+\varepsilon_2+\varepsilon_3\in\{3,1,-1,-3\}$
with multiplicities $1,3,3,1$. Reordering $\omega_1\omega_2\omega_3=e_1e_4e_2e_5e_3e_6$ into $e_1\cdots e_6$ costs three
transpositions, so $\omega=-\omega_1\omega_2\omega_3=-\varepsilon_1\varepsilon_2\varepsilon_3$ on that line. Hence $\tfrac12(\varpi+\omega)$ takes the
values $1,1,-1,-1$ on the four strata: it is the sign of $\varpi$, and equals the polynomial $(13\varpi-\varpi^3)/12$.
The operator
\begin{equation}
\label{eq:Piclosed}
\Pi:=\tfrac12\id-\tfrac14(\varpi+\omega)
\end{equation}
is therefore the orthogonal projector of rank four onto the strata with $\varpi<0$. We call it the paired-mode projector; it
is a natural candidate for an observation projector at one constituent and we test it against the photon in
\cref{sec:photon}. It is not Lorentz invariant: of the six interior Lorentz generators $e_1e_4,e_1e_5,e_1e_6,e_4e_5,e_4e_6,e_5e_6$ only
$e_1e_4$ commutes with $\Pi$ (the commutators have ranks $0,4,4,4,4,4$), and $[\Pi,J]$ has rank four. Its stabiliser in $\so{3,3}$ is
nine-dimensional, a count of linear combinations, not of commuting basis bivectors.
\end{example}

% --- source: tex/s02_arena#5 (The interior four-plane and the surplus plane)
\subsection{The interior four-plane and the surplus plane}
\label{sec:interior}

The physical reading of the arena singles out a Lorentzian four-plane inside $V$. We take
\[
E=\langle e_1,e_4,e_5,e_6\rangle,\qquad \gamma^0=e_1,\quad \gamma^a=e_{a+3}\ (a=1,2,3),
\]
of signature $(1,3)$, and call the orthogonal complement $U=\langle e_2,e_3\rangle$ the surplus plane. The interior Lorentz
algebra $\lint$ is spanned by the six bivectors of $E$, with rotations $S_1=\tfrac12\gamma^2\gamma^3$ (cyclically) and boosts
$T_a=\tfrac12\gamma^0\gamma^a$. Its spin Casimir on $S$ is $-\sum_aS_a^2=\tfrac34\id$, and a rotation through $2\pi$ is $-\id$
on $S$ while it is $+\id$ on vectors (\cref{sec:spin}): the same rotors give spin one-half to spinors and spin one to vectors.

\begin{proposition}[The interior charge is unique]
\label{prop:6plus8plus1}
As $\lint$-modules, $\Lambda^2V=\Lambda^2E\oplus(E\otimes U)\oplus\Lambda^2U$, of dimensions $6+8+1$, is the adjoint
representation, two copies of the vector representation, and the trivial representation. The centraliser of $\lint$ in
$\so{3,3}$ is therefore the line spanned by
\[
J:=e_2e_3 ,
\]
which satisfies $J^2=-\id$, $J^{\!\top}=-J$, and commutes with each of the four gamma matrices.
\end{proposition}

\begin{proof}
The adjoint and vector representations of $\so{1,3}$ contain no invariant, so only $\Lambda^2U$ survives. That $J$ commutes
with $e_1,e_4,e_5,e_6$ is \eqref{eq:bladecomm} with disjoint index sets and $|I||J|=2$.
\end{proof}

\begin{figure}[t]
\centering
\begin{tikzpicture}[>=Stealth]
\coordinate (t1) at (0,2.6); \coordinate (t2) at (0,1.3); \coordinate (t3) at (0,0);
\coordinate (s1) at (5,2.6); \coordinate (s2) at (5,1.3); \coordinate (s3) at (5,0);
% interior four-plane (L-shaped region)
\fill[gcol!8,rounded corners=8pt] (-1.85,3.05) -- (7.0,3.05) -- (7.0,-0.5) -- (4.45,-0.5) -- (4.45,2.0) -- (-1.85,2.0) -- cycle;
\draw[gcol,very thick,rounded corners=8pt] (-1.85,3.05) -- (7.0,3.05) -- (7.0,-0.5) -- (4.45,-0.5) -- (4.45,2.0) -- (-1.85,2.0) -- cycle;
% surplus plane
\draw[ocol,very thick,dashed,rounded corners=6pt] (-0.75,-0.5) rectangle (0.55,1.75);
% pairing arrows
\draw[->,gray,thick,shorten <=4pt,shorten >=4pt] (t1) -- node[above] {\small $e_i\mapsto e_{i+3}$} (s1);
\draw[->,gray,thick,shorten <=4pt,shorten >=4pt] (t2) -- (s2);
\draw[->,gray,thick,shorten <=4pt,shorten >=4pt] (t3) -- (s3);
\foreach \p in {t1,t2,t3}{\fill[tcol] (\p) circle (2.4pt);}
\foreach \p in {s1,s2,s3}{\fill[scol] (\p) circle (2.4pt);}
\node[left=3pt] at (t1) {$e_1=\gamma^0$}; \node[left=3pt] at (t2) {$e_2$}; \node[left=3pt] at (t3) {$e_3$};
\node[right=3pt] at (s1) {$e_4=\gamma^1$}; \node[right=3pt] at (s2) {$e_5=\gamma^2$}; \node[right=3pt] at (s3) {$e_6=\gamma^3$};
\node[tcol,anchor=south] at (0,3.15) {\small $V^+$ (times): $e_i^2=+\id$};
\node[scol,anchor=south] at (5,3.15) {\small $V^-$ (lengths): $e_{i+3}^2=-\id$};
\node[ocol,anchor=east,align=right] at (-0.95,0.6) {\small surplus plane $U$,\\ \small $J=e_2e_3$};
\node[gcol,anchor=north,align=center] at (5.7,-0.6) {\small interior four-plane\\ \small $E=\langle e_1,e_4,e_5,e_6\rangle$, signature $(1,3)$};
\end{tikzpicture}
\caption{The arena. The positive triple $V^+$ (times) and the negative triple $V^-$ (lengths) are paired by $e_i\mapsto e_{i+3}$;
the pairing has two square roots, the grading $\cK$ and the binding operator $G$ (\cref{prop:twoK}). The interior four-plane
$E$, on which the Dirac sector is built, contains the clock $e_1$ and the whole spatial triple; its complement is the surplus
plane $U=\langle e_2,e_3\rangle$, whose rotation generator $J$ is the only bivector commuting with the interior Lorentz algebra
(\cref{prop:6plus8plus1}).}
\label{fig:arena}
\end{figure}

\Cref{prop:6plus8plus1} is the first instance of a pattern that recurs throughout: a structure is fixed once the right
ambient algebra and the right conditions are named, and it ceases to be fixed when either is relaxed. Here the ambient
space is the span of the bivectors and the condition is commutation with $\lint$. In \cref{sec:dirac} we enlarge the ambient
algebra to all of $\End_\R(S)$ and find that a different condition, commutation with the internal bivectors, is what keeps the
answer unique.

\section{The spinor module as forms}
\label{sec:w:thespinormoduleasforms}

% --- source: two_readings/r3_octants#2 (The octants and the octant group)
\subsubsection{The octants and the octant group}
\label{r3:sec:octants}

\begin{proposition}[Octants and Fock grading]
\label{r3:prop:octants}
\begin{enumerate}[label=(\roman*),leftmargin=2.2em]
\item The pole operators are $\omega_i|b\rangle=(-1)^{b_i}|b\rangle$ on the basis $|b_1b_2b_3\rangle$ of $S$; their joint eigenlines, the
\emph{octants}, are the eight basis lines, labelled by the sign triples $\sigma=(\sigma_1,\sigma_2,\sigma_3)$, $\sigma_i=(-1)^{b_i}$.
\item $a^\pm_i:=\tfrac12(e_i\pm e_{i+3})$ satisfy $\{a^+_i,a^+_j\}=\{a^-_i,a^-_j\}=0$, $\{a^+_i,a^-_j\}=\delta_{ij}$, and
$\omega_i=\id-2a^+_ia^-_i$. Hence $N:=\sum_ia^+_ia^-_i$ is diagonal and counts the minus signs of $\sigma$.
\item $|\mathrm{vac}\rangle:=|000\rangle$, of signs $(+,+,+)$, spans the common kernel of the $a^-_i$ (the Fock vacuum, $N=0$), and
$a^+_{i_1}\cdots a^+_{i_r}|\mathrm{vac}\rangle=|b_I\rangle$ for $i_1<\dots<i_r$, where $b_I$ has ones exactly at $I$. The filled pole is
$|\mathrm{full}\rangle:=|111\rangle$, of signs $(-,-,-)$ and $N=3$.
\item The octant $\sigma$ is the pure spinor line of the maximal isotropic plane $N_\sigma:=\operatorname{span}\{e_i-\sigma_ie_{i+3}\}$: the
vectors $v\in V$ with $v|b\rangle=0$ are exactly those of $N_\sigma$.
\item $\Gamma=\omega=-\omega_1\omega_2\omega_3$, so $(-1)^N=\omega_1\omega_2\omega_3=-\Gamma$; the element $\Gamma$ is, up to the sign fixed by the choice $i=\pm J$, the chirality $\gamma_5$ of \cref{sec:discrete}. On the octant with $w=3-N$ plus signs, $\Gamma=(-1)^w$.
\end{enumerate}
\end{proposition}

\begin{proof}
(i) is in \cref{prop:grading} and \cref{app:realisation}: $e_i|b\rangle=(-1)^{b_1+\dots+b_{i-1}}|b\oplus\delta_i\rangle$ and
$e_{i+3}|b\rangle=(-1)^{b_1+\dots+b_i}|b\oplus\delta_i\rangle$. (ii) The Clifford relations give $\{e_i\pm e_{i+3},e_j\pm e_{j+3}\}=2(\delta_{ij}-\delta_{ij})=0$
and $\{e_i+e_{i+3},e_j-e_{j+3}\}=4\delta_{ij}$; then $\omega_i=(a^+_i+a^-_i)(a^+_i-a^-_i)=a^-_ia^+_i-a^+_ia^-_i=\id-2a^+_ia^-_i$. (iii)/(iv) By the two
formulas of (i), $(e_i-s\,e_{i+3})|b\rangle=(-1)^{b_1+\dots+b_{i-1}}(1-s(-1)^{b_i})|b\oplus\delta_i\rangle$, which vanishes iff $s=\sigma_i$. So $N_\sigma$
annihilates the octant $\sigma$; an annihilator is isotropic, since $2\eta(u,v)\phi=(uv+vu)\phi=0$, hence of dimension at most three, and it equals
$N_\sigma$. For $\sigma=(+,+,+)$ this is the statement $a^-_i|000\rangle=0$. The sign of $a^+_{i_1}\cdots a^+_{i_r}|\mathrm{vac}\rangle$ is a product of
factors $(-1)^{b_1+\dots+b_{i-1}}$ taken while all bits before $i$ are still zero. (v) Reordering $e_1e_4e_2e_5e_3e_6$ into $e_1\cdots e_6$
costs three transpositions (\cref{ex:pairedmode}).
\end{proof}

All five items are restated from \cref{sec:arena} or fixed here (r3-1, r3-2, r3-3). Item (v) gives $\Gamma=-\omega_1\omega_2\omega_3$ with eigenvalue $(-1)^w$, where the temporal weight $w=3-N$ counts plus signs. The element $\omega=e_1\cdots e_6$ of \eqref{eq:tower}, the $\Gamma$ of (v) and of \cref{thm:i}, and the volume element used in \cref{ex:pairedmode} are one and the same.

\begin{proposition}[The octant group]
\label{r3:prop:octgroup}
For each axis $a$, the automorphisms of $\Cl(3,3)$ extending the reflection $R_a:e_a\mapsto-e_a$, $e_b\mapsto e_b$ $(b\neq a)$ are the
conjugations by the multiples of $g_a:=e_a\omega$. Each $g_a$ is a signed permutation matrix; it conjugates $\omega_i$ to $-\omega_i$ for
$a\in\{i,i+3\}$ and fixes the other two, so it carries the octant $\sigma$ to the octant with $\sigma_i$ reversed. The reflections of a
temporal axis and of its spatial partner act alike on the octants, and $g_ig_{i+3}=\pm\omega_i$ fixes every octant. The lifts $g_1,g_2,g_3$
satisfy $g_i^2=-\id$, anticommute pairwise and generate a group of order $16$ whose centre $\{\pm\id,\pm g_1g_2g_3\}$ contains
$g_1g_2g_3=e_4e_5e_6$, the element mapping each octant to its antipode; the kernel of its action on the eight octants is $\{\pm\id\}$, and
it acts through $\Z_2^3$, simply transitively.
\end{proposition}

\begin{proof}
$\Cl(3,3)\cong M_8(\R)$ is central simple, so every automorphism is inner and its implementing element is unique up to a real factor
(Skolem--Noether). Since $\omega$ anticommutes with every generator, $g_ae_bg_a^{-1}=-e_ae_be_a^{-1}$, which is $-e_b$ for $b=a$ and $e_b$
otherwise. Hence $g_a\omega_ig_a^{-1}=R_a(e_i)R_a(e_{i+3})$. Further $g_i^2=-e_i^2\omega^2=-\id$, and $g_ig_j=-e_ie_j$ anticommutes with
$g_jg_i=-e_je_i$. Moving $g_i$ past the other two lifts costs two signs, so $g_1g_2g_3=-e_1e_2e_3\,\omega=e_4e_5e_6$ commutes with every $g_i$; it is
not $\pm\id$, and it reverses all three signs. The images of $|\mathrm{vac}\rangle$ under the eight sign classes are the eight octants
(r3-4, r3-5).
\end{proof}

Fixed. One may also define the octant group as the reflections of the temporal $\R^3$ acting on its eight \emph{sign sectors}
$\{x\in V^+:\sigma_ix_i>0\}$, and identify them with the spinor octants. Both sets are $\Z_2^3$-torsors, so equivariant bijections
between them exist, eight of them, one for each choice of the octant assigned to the positive sector. One of them is singled out by an operator of the clock analysis (\cref{sec:clock}):

\begin{proposition}[Clock octant and vacuum]
\label{r3:prop:clock}
For a dimensionless $x\in V^+$ the conversion operator of \cref{thm:spectrumBC} with $C=\operatorname{diag}(x)$ is
$B_x=\sum_ix_ie_ie_{i+3}=\sum_ix_i\omega_i$. It has eigenvalue $\sum_i\sigma_ix_i$ on the octant $\sigma$; if no $x_i$ vanishes, its unique top
eigenline is the octant $\operatorname{sign}(x)$; and $g_iB_xg_i^{-1}=B_{R_ix}$. The Fisher--Rao octant of clocks of \cref{sec:clocksimplex},
the closed positive octant $x_i=2\sqrt{w_i}\ge0$ of the sphere of radius two, therefore corresponds to the Fock vacuum: for $x$ in its
interior the top eigenline of $B_x$ is $|\mathrm{vac}\rangle$. At the symmetric point the radius-two coordinates $x_i=2/\sqrt3$ give
$B_x=2\varpi/\sqrt3$, twice the unit clock $B_{\widehat W}=\varpi/\sqrt3$, whose top eigenvalue $\sqrt3$ lies on $|\mathrm{vac}\rangle$; the top
eigenspaces agree. The weights $w_i=x_i^2/4$ do not see the signs of the amplitudes, and the clock simplex is the quotient of the clock sphere by
the octant group. The quotient has boundary stabilisers: if $k$ coordinates of $x$ vanish, the top eigenspace of $B_x$ has dimension $2^k$
($8$ at $x=0$), and a point on such a face of the simplex has $2^{3-k}$ preimages. The map is not an eight-sheeted covering.
\end{proposition}

\begin{proof}
$B_x$ is diagonal by \cref{r3:prop:octants}(i), and $\sum_i\sigma_ix_i<\sum_i|x_i|$ unless $\sigma=\operatorname{sign}(x)$. Equivariance follows
from \cref{r3:prop:octgroup} (r3-6, with twenty random rational $x$).
\end{proof}

The spectral statements are fixed. The identification of the clock octant with the vacuum pole is constructed: it uses the order in
$\omega_i=e_ie_{i+3}$ and the choice of the top rather than the bottom eigenline, and the opposite choice assigns the filled pole. The two
objects are different in kind, a two-dimensional region of the sphere of clock directions in $V^+$ and a line of the spinor module; read
physically, the correspondence is interpretive.

% --- source: two_readings/r3_octants#3 (The exterior-algebra structure)
\subsubsection{The exterior-algebra structure}
\label{r3:sec:exterior}

\begin{proposition}[The spinor module as an exterior algebra]
\label{r3:prop:fock}
$L:=\operatorname{span}\{e_i+e_{i+3}\}$ and $L':=\operatorname{span}\{e_i-e_{i+3}\}$ are maximal isotropic, transverse, and dual to each other
under $\eta$. The map $\Lambda^r(L)\to S$, $\ell_1\wedge\dots\wedge\ell_r\mapsto2^{-r}\ell_1\cdots\ell_r|\mathrm{vac}\rangle$, is an isomorphism
$\Lambda^*(L)\cong S$ carrying the degree to $N$, with $\dim\Lambda^r=1,3,3,1$ for $r=N=0,1,2,3$: the weights of the octants. Writing
$X=X_L+X_{L'}$, a vector acts as exterior multiplication by $2X_L$ plus contraction with the form $\eta(X_{L'},\cdot)$ on $L$,
\begin{equation}
\label{r3:eq:Xfock}
X=\sum_i\bigl[(B_i+E_i)\,a^+_i+(B_i-E_i)\,a^-_i\bigr].
\end{equation}
Moreover $a\mapsto a|\mathrm{vac}\rangle$ is an isomorphism from the temporal Clifford algebra $\Cl(V^+)=\langle e_1,e_2,e_3\rangle\cong\Cl(3,0)$ onto
$S$, and $e_I|\mathrm{vac}\rangle=a^+_I|\mathrm{vac}\rangle$: the octant weight $N$ is the Clifford grade of the temporal blade $e_I$.
\end{proposition}

\begin{proof}
$\eta(e_i\pm e_{i+3},e_j\pm e_{j+3})=\delta_{ij}-\delta_{ij}=0$, while $\eta(e_i+e_{i+3},e_j-e_{j+3})=2\delta_{ij}$. Elements of $L$ anticommute, so the
map is well defined, and it is bijective by \cref{r3:prop:octants}(iii). Under it $a^+_i$ is $\ell_i\wedge$ with $\ell_i=e_i+e_{i+3}=2a^+_i$, and
$\ell'_i=e_i-e_{i+3}=2a^-_i$ is the contraction with $\eta(\ell'_i,\cdot)$, since $\eta(\ell'_i,\ell_j)=2\delta_{ij}$. \Cref{r3:eq:Xfock} is
$e_i=a^+_i+a^-_i$, $e_{i+3}=a^+_i-a^-_i$. Finally
$e_j=a^+_j+a^-_j$ and $a^-_j$ annihilates $a^+_I|\mathrm{vac}\rangle$ for $j\notin I$, which gives $e_I|\mathrm{vac}\rangle=a^+_I|\mathrm{vac}\rangle$ by induction (r3-7).
\end{proof}

Inherited: this is Chevalley's realisation of the spinors of a space $L\oplus L'$ with $L'\cong L^*$ on $\Lambda^*(L)$
\cite{Chevalley54,LawsonMichelsohn}, and the pure spinor lines of \cref{r3:prop:octants}(iv) are the lines of Cartan and Chevalley (see Section~2
of \cite{Gualtieri03}). The last statement is the classical identification of the spinors of $\Cl(3,3)\cong M_8(\R)$ restricted to
$\Cl(3,0)$ with the Pauli algebra itself. The same realisation of the spinors of $\mathrm{SO}(n,n)$ by the forms on a null $n$-plane is the
starting point of Krasnov \cite{Krasnov18}, who reads the forms as fermion fields; for $\mathrm{SO}(3,3)$ he finds the commutant
$\mathrm{SO}(2)$ of the Lorentz group and a single charged two-component Weyl fermion. Here the forms carry the field strength,
$\Phi_{\rm oct}(X)=\bE\cdot f+\bB\cdot g$ (\cref{r3:thm:phi}), and \cref{r3:thm:noeven} shows that this field reading is not the spin representation.
The module $\Lambda^*(L)$ is shared; the physical identification is not.

\begin{proposition}[The commutant of the grading root]
\label{r3:prop:gl3}
The elements of $\so{3,3}$ commuting with the grading root $\cK=\tfrac12\ad_{\varpi}|_V$ (\cref{prop:twoK}) form a subalgebra of dimension nine, the adjoint image of the bivectors
$E_{ij}:=a^+_ia^-_j-\tfrac12\delta_{ij}\id$ (so $E_{ij}=a^+_ia^-_j$ for $i\neq j$ and $E_{ii}=-\tfrac12\omega_i$). They commute with $N$; $\ad E_{ij}$ acts
on $L$ by the matrix unit $e_{ij}$ and on $L'$ by $-e_{ji}$, so the subalgebra is $\mathfrak{gl}(L)\cong\mathfrak{gl}(3,\R)$; and the spin lift of
$A\in\mathfrak{gl}(3)$ acts on $\Lambda^r(L)$ as $d\Lambda^r(A)-\tfrac12\tr A$, where $d\Lambda^r(A)$ is the derived action, extending $A$ to
$\Lambda^r(L)$ as a derivation (for $A=\id$ the eigenvalue on degree $r$ is $r-\tfrac32$). On $GL^+(L)$ the integrated action is
$\rho(g)=(\det g)^{-1/2}\Lambda^*(g)$. As a $\mathfrak{gl}(3)$-module, $S\cong\Lambda^*(L)\otimes(\det L)^{-1/2}$.
\end{proposition}

\begin{proof}
$\cK$ has the eigenspaces $L$ (eigenvalue $-1$) and $L'$ ($+1$), so an element commuting with it preserves both; $\eta$-skewness then
forces $A|_{L'}=-(A|_L)^*$ under $L'\cong L^*$, and every $A|_L\in\mathfrak{gl}(L)$ occurs: dimension nine. Since $\cK=\tfrac12\ad_{\varpi}$ and
$\varpi=3\,\id-2N$, an element $\ad\beta$ commutes with $\cK$ iff $[N,\beta]=0$. A bivector is quadratic in the $a^\pm$; the number-conserving
ones are the $a^+_ia^-_j$ up to scalars, and $E_{ij}$ removes the scalar part. By the CAR $[a^+_ia^-_j,a^+_k]=\delta_{jk}a^+_i$ and
$[a^+_ia^-_j,a^-_k]=-\delta_{ik}a^-_j$, and $a^+_ia^-_j$ acts on $a^+_K|\mathrm{vac}\rangle$ as the derivation replacing $a^+_j$ by $a^+_i$ (r3-8).
\end{proof}

The statement is standard (Example~2.12 of \cite{Gualtieri03} gives $\Lambda^*V^*\otimes(\det V)^{1/2}$ with $V^*$ in the role of $L$); it is fixed
here in the realisation. It shows that $\mathfrak{gl}(3)$ does not act on $\Lambda^r$ by the bare derived action: the density factor
$-\tfrac12\tr A$, the same on all degrees, is part of the spin representation.

\subsection{The trace form splits by the Cartan involution}
\label{sec:cartansplit}

\begin{theorem}[The trace form splits by the Cartan involution]
\label{thm:ghost}
Let $P=P^{\!\top}=P^2$, let $B$ be a bivector and $\Delta=[B,P]$. Then $\Delta^{\!\top}=PB^{\!\top}-B^{\!\top}P$, and
\begin{enumerate}[label=(\roman*),leftmargin=2.2em]
\item if $B$ is compact, $B^{\!\top}=-B$, then $\Delta$ is symmetric and $\Tr(\Delta^2)=\Tr(\Delta^{\!\top}\Delta)>0$ unless $\Delta=0$;
\item if $B$ is a boost, $B^{\!\top}=B$, then $\Delta$ is antisymmetric and $\Tr(\Delta^2)=-\Tr(\Delta^{\!\top}\Delta)<0$ unless $\Delta=0$.
\end{enumerate}
Symmetric and antisymmetric matrices are trace-orthogonal, so on an orbit the trace form is an orthogonal sum of a positive
definite piece on the directions generated by compact bivectors and a negative definite piece on those generated by boosts. Only
the former are tangent to the manifold of orthogonal projectors.
\end{theorem}

\begin{proof}
$\Delta^{\!\top}=(BP-PB)^{\!\top}=PB^{\!\top}-B^{\!\top}P$, which is $\Delta$ when $B^{\!\top}=-B$ and $-\Delta$ when $B^{\!\top}=B$. The trace of the
product of a symmetric and an antisymmetric matrix vanishes. A curve of orthogonal projectors has symmetric velocity, so the
antisymmetric directions leave that manifold.
\end{proof}

\Cref{thm:sig} and \cref{thm:ghost} classify bivectors relative to the fixed Cartan decomposition, compact meaning antisymmetric for the
chosen positive form. A general element $X=X_a+X_s$ of $\so{3,3}$, split into antisymmetric and symmetric parts, gives a tangent to the
manifold of orthogonal projectors exactly when $[X_s,P]=0$, so generators are to be counted modulo the stabiliser of $P$. The theorem fixes the
sign of the trace form on an orbit; it fixes the sign of a kinetic term only for actions built on this trace metric, not that of an arbitrary
action, spacetime contraction, constraint or propagator residue.

